\centerline{\bf \Huge Ортоцентр}

\begin{flushright}
        \scriptsize{Эпизод 9. Действовать синхронно}
\end{flushright}

\problem{}
Предположим, что две вершины треугольника и его описанная окружность зафиксированы, а третья вершина движется по этой описанной окружности. По какой траектории будет двигаться ортоцентр треугольника? (Не забудьте, что вершина может двигаться по двум разным дугам!)

\problem{}
Обозначим за $O$ и $H$ центр описанной около треугольника $A B C$ окружности и его ортоцентр. Докажите, что
(a) $\angle B A H=\angle C A O$
(b) Расстояние от $O$ до $B C$ равно половине отрезка $A H$.

\problem{}
В остроугольном треугольнике $A B C$ угол при вершине $A$ равен $60^{\circ}$. Известно, что $B C=1$. Найдите длину отрезка $B_1 C_1$, если $B_1$ и $C_1$ это основания высот треугольника, опущенных из вершин $B$ и $C$.

\problem{}
Пусть $H$ - ортоцентр остроугольного треугольника $A B C$. Точка $A^{\prime}$ выбирается на описанной около этого треугольника окружности $\Omega$ так, что $A A^{\prime}$ - диаметр $\Omega$. Докажите, что четырехугольник $B H C A^{\prime}$ - параллелограмм.

\problem{}
Треугольник $A B C$ с ортоцентром $H$ вписан в окружность $\omega$. Окружность, построенная на отрезке $A H$ как на диаметре, пересекает $\omega$ в точке $S \neq A$. Докажите, что прямая $S H$ делит отрезок $B C$ пополам.

\problem{}
Четырехугольник $A B C D$ вписанный. Докажите, что ортоцентры треугольников $A B C, B C D$, $C D A$ и $D A B$ образуют четырехугольник, равный исходному.

\problem{}
Точка Анти-Штейнера. Прямая $\ell$ проходит через ортоцентр треугольника $A B C$. Докажите, что три прямые, симметричные прямой $\ell$ относительно сторон треугольника, пересекаются в одной точке.